% Section 2: Tensors of type $(\ell,k)$} | Part: 1 | Depends on: see docs/REWRITE_PLAN.md | Status: todo | Original lines: 654-704
%!TEX root = ../../../main.tex
\subsection{Tensors of type $(\ell,k)$}\label{sec:tensortypes}
\begin{formal}[{tensors}]
A \emph{mixed tensor, $k$-covariant and $\ell$-contravariant} on $E$ is a multilinear map
\[
\tau:\underbrace{E\times\dots\times E}_{k}\times\underbrace{E^*\times\dots\times E^*}_{\ell}\to\R .
\]
The space of such tensors is $\TT{\ell}{k}(E)\cong \underbrace{E^*\otimes\dots\otimes E^*}_{k}\otimes\underbrace{E\otimes\dots\otimes E}_{\ell}$. Conventionally $\TT{0}{0}(E)=\R$ (scalars). Special cases:
\begin{center}
\begin{tabular}{@{}lll@{}}
\toprule
Space & Meaning & Components\\ \midrule
$\TT{0}{0}=\R$ & scalar & $\tau$\\
$\TT{0}{1}=E^*$ & covector & $f_i$\\
$\TT{1}{0}=(E^*)^*\cong E$ & vector & $v^i$\\
$\TT{0}{2}$ & bilinear form (\eg metric) & $g_{ij}$\\
$\TT{1}{1}$ & linear map $E\to E$ (via $A(v,\alpha)=\alpha(Av)$) & $A^i{}_j$\\
$\TT{1}{2}$ & (\eg torsion) & $T^i{}_{jk}$\\
$\TT{1}{3}$ & (\eg Riemann curvature) & $R^l{}_{ijk}$\\ \bottomrule
\end{tabular}
\end{center}
\end{formal}

\begin{intuition}[{a tensor is a coordinate-free machine; its components are a table}]
A tensor is a geometric object like a vector: it exists independently of any basis. It is a machine with $k$ ``vector hoppers'' and $\ell$ ``covector hoppers'' that returns a number. The components are just the table of outputs on basis elements:
\[
\tau^{i_1\dots i_\ell}{}_{j_1\dots j_k}=\tau\big(e_{j_1},\dots,e_{j_k},e^{i_1},\dots,e^{i_\ell}\big).
\]
So $\tau=\tau^{i_1\dots i_\ell}{}_{j_1\dots j_k}\;e^{j_1}\otimes\dots\otimes e^{j_k}\otimes e_{i_1}\otimes\dots\otimes e_{i_\ell}$.
\end{intuition}

\begin{formal}[{transformation rule of components}]
If $e'_i=A^j{}_ie_j$, the components of a tensor in the new basis are
\begin{equation}\label{eq:tensortransf}
\tau'^{\,i_1\dots i_\ell}{}_{j_1\dots j_k}=(A^{-1})^{i_1}{}_{a_1}\cdots(A^{-1})^{i_\ell}{}_{a_\ell}\;A^{b_1}{}_{j_1}\cdots A^{b_k}{}_{j_k}\;\tau^{a_1\dots a_\ell}{}_{b_1\dots b_k}.
\end{equation}
Each upper index is transformed with $A^{-1}$, each lower index with $A$. A physicist's definition: \emph{a tensor is an array of numbers that transforms like this under changes of basis}. (In the notes this appears with $B=A$.)
\end{formal}

\begin{example}[{three ``matrices'', three different transformation laws}]
In matrix notation, with $M$ the component matrix and $A$ the change of basis:
\[
\TT{0}{2}\ (g_{ij}):\ M'=A^tMA,\qquad
\TT{1}{1}\ (A^i{}_j):\ M'=A^{-1}MA,\qquad
\TT{2}{0}\ (T^{ij}):\ M'=A^{-1}M A^{-t}.
\]
\emph{Test.} Dot product on $\R^2$, $g=\mathrm{diag}(1,1)$, and $A=\begin{pmatrix}1&-1\\1&1\end{pmatrix}$ from Section~\ref{sec:vs}: $g'=A^tA=\begin{pmatrix}2&0\\0&2\end{pmatrix}$: the new basis vectors are still orthogonal but have length $\sqrt2$. \emph{Test 2.} The trace of a $(1,1)$ tensor is invariant, $\tr(A^{-1}MA)=\tr M$, but $\tr(A^tMA)\neq\tr M$ in general: ``trace of a bilinear form'' is not a meaningful invariant without a metric.
\end{example}
